%# -*- coding:utf-8 -*-
% !TEX program = xelatex
%% start of file `template_en.tex'.
%% Copyright 2006-1008 Xavier Danaux (xdanaux@gmail.com).
%
% This work may be distributed and/or modified under the
% conditions of the LaTeX Project Public License version 1.3c,
% available at http://www.latex-project.org/lppl/.


\documentclass[11pt,a4paper]{moderncv}

\usepackage{fontspec,xunicode}
\usepackage[slantfont,boldfont]{xeCJK}
\usepackage{xcolor}  % replace by the encoding you are using

\usepackage{ifplatform} % 判断平台


% \setmainfont{Tahoma}
\setmainfont{DejaVu Serif}  % 缺省英文字体.serif是有衬线字体sans serif无衬线字体

\ifwindows
  % Windows 下使用的代码
  \setCJKmainfont[ItalicFont={AR PL KaitiM GB}, BoldFont={Droid Sans Fallback}]{AR PL SungtiL GB}  % 衬线字体 缺省中文字体为
  \setCJKsansfont{AR PL SungtiL GB}
  \setmonofont{FangSong}                 % 等宽字体: FangSong, Consolas
\fi

\ifmacosx
  % macOS 下使用的代码
  \setCJKmainfont[ItalicFont={AR PL KaitiM GB}, BoldFont={Droid Sans Fallback}]{AR PL SungtiL GB}  % 衬线字体 缺省中文字体为
  \setCJKsansfont{STSong}
  \setCJKmonofont{STFangsong}  % 中文等宽字体
\fi


%-----------------------xeCJK下设置中文字体------------------------------%
\setCJKfamilyfont{song}{AR PL SungtiL GB}  % 宋体 song
\newcommand{\song}{\CJKfamily{song}}
\setCJKfamilyfont{fs}{FangSong_GB2312}  % 仿宋2312 fs
\newcommand{\fs}{\CJKfamily{fs}}
\setCJKfamilyfont{yh}{Microsoft YaHei}  % 微软雅黑 yh
\newcommand{\yh}{\CJKfamily{yh}}
\setCJKfamilyfont{hei}{SimHei}  % 黑体  hei
\newcommand{\hei}{\CJKfamily{hei}}
\setCJKfamilyfont{hwxh}{STXihei}  % 华文细黑  hwxh
\newcommand{\hwxh}{\CJKfamily{hwxh}}
\setCJKfamilyfont{asong}{Adobe Song Std}  % Adobe 宋体  asong
\newcommand{\asong}{\CJKfamily{asong}}
\setCJKfamilyfont{ahei}{Adobe Heiti Std}  % Adobe 黑体  ahei
\newcommand{\ahei}{\CJKfamily{ahei}}
\setCJKfamilyfont{akai}{Adobe Kaiti Std}  % Adobe 楷体  akai
\newcommand{\akai}{\CJKfamily{akai}}


%------------------------------设置字体大小------------------------%
\newcommand{\chuhao}{\fontsize{42pt}{\baselineskip}\selectfont}  % 初号
\newcommand{\xiaochuhao}{\fontsize{36pt}{\baselineskip}\selectfont}  % 小初号
\newcommand{\yihao}{\fontsize{28pt}{\baselineskip}\selectfont}  % 一号
\newcommand{\erhao}{\fontsize{21pt}{\baselineskip}\selectfont}  % 二号
\newcommand{\xiaoerhao}{\fontsize{18pt}{\baselineskip}\selectfont}  % 小二号
\newcommand{\sanhao}{\fontsize{15.75pt}{\baselineskip}\selectfont}  % 三号
\newcommand{\sihao}{\fontsize{14pt}{\baselineskip}\selectfont}  % 四号
\newcommand{\xiaosihao}{\fontsize{12pt}{\baselineskip}\selectfont}  % 小四号
\newcommand{\wuhao}{\fontsize{10.5pt}{\baselineskip}\selectfont}  % 五号
\newcommand{\subwuhao}{\fontsize{10pt}{\baselineskip}\selectfont}  % 次五号
\newcommand{\xiaowuhao}{\fontsize{9pt}{\baselineskip}\selectfont}  % 小五号
\newcommand{\liuhao}{\fontsize{7.875pt}{\baselineskip}\selectfont}  % 六号
\newcommand{\qihao}{\fontsize{5.25pt}{\baselineskip}\selectfont}  % 七号


% \usepackage{fontawesome}
% \setCJKmainfont[BoldFont={WenQuanYi Micro Hei/Bold}]{WenQuanYi Micro Hei}
% \defaultfontfeatures{Mapping=tex-text}
% \XeTeXlinebreaklocale "zh"
% \XeTeXlinebreakskip = 0pt plus 1pt minus 0.1pt
% moderncv themes
\moderncvstyle{classic}
\moderncvcolor{blue}
% \moderncvtheme[green]{classic}  % idem
% \moderncvtheme[blue,roman]{hht}
% character encoding



% adjust the page margins
\usepackage[scale=0.9]{geometry}
\setlength{\hintscolumnwidth}{3cm}  % if you want to change the width of the column with the dates
% \AtBeginDocument{\setlength{\maketitlenamewidth}{6cm}}  % only for the classic theme, if you want to change the width of your name placeholder (to leave more space for your address details
\AtBeginDocument{\recomputelengths}  % required when changes are made to page layout lengths

% personal data
\firstname{张}
\familyname{奥恒}
\title{ZhangAoheng}  % optional, remove the line if not wanted
% \address{杭州}{}  % optional, remove the line if not wanted
\address{安徽大学计算机学院}{}  % optional, remove the line if not wanted
\mobile{18539197010}  % optional, remove the line if not wanted
% \fax{fax (optional)}  % optional, remove the line if not wanted
\email{2389009059@qq.com}  % optional, remove the line if not wanted
\extrainfo{%
  WeChat: ZhangaooooooHeng \\
 }

\photo[64pt]{avatar.png}  % '64pt' is the height the picture must be resized to and 'picture' is the name of the picture file; optional, remove the line if not wanted
% \quote{China\TeX 您的LaTeX乐园，TeX\&\LaTeX 王国}  % optional, remove the line if not wante

% \nopagenumbers{}  % uncomment to suppress automatic page numbering for CVs longer than one page


%----------------------------------------------------------------------------------
%            content
%----------------------------------------------------------------------------------
\begin{document}
\maketitle
\vspace*{-14mm}

\section{教育经历}
\cventry{11.09-15.06}{本科}{河南师范大学}{计算机科学与技术}{}{}  % arguments 3 to 6 are optional
\cvlistitem{全国大学生创新创业大赛二等奖}
\cvlistitem{国家奖学金/学生会主席/优秀毕业生}
\cventry{15.09-18.06}{硕士}{安徽大学}{计算机科学与技术}{}{}  % arguments 3 to 6 are optional
\cvlistitem{通过 CET-6，具备良好的英文文献阅读能力}

\section{项目}
\subsection{科研项目}
\cvline{基于道路标识的车辆定位方法研究}{\textbf{项目简介：}提出了一种基于路牌语义地标的视觉 SLAM 回环检测方法，因天气变化，而造成的室外环境的变化，传统视觉SLAM无法进行的回环检测，核心创新是利用车辆靠近路牌时其在图像中的膨胀效应（Looming），结合IMU提供的真实尺度轨迹，估算出路牌的绝对深度，突破了纯单目视觉无法获取真实尺度的难题。在恶劣天气下仍能实现稳定的回环闭合与轨迹校正。主要工作：}
\cvline{}{\textbf{前端视觉检测：}使用YOLO网络进行道路标识牌区域检测，对道路表示牌外轮廓进行线检测，确立道路表示牌中心点。对道路表示牌内部汉字进行线检测，为后端位姿估计提供数据。}
\cvline{}{\textbf{后端深度估计：}结合ORB-SLAM3算法所计算出的位姿数据，计算FOE（膨胀焦点），计算路牌中心到FOE的径向距离（膨胀量），利用径向膨胀量与帧间位移的比例关系解出路牌的绝对深度（Looming测距）。}
\cvline{}{\textbf{后端位姿计算：}对跨序列的同一路牌进行 LoFTR 特征匹配，求解单应性矩阵并分解得到多组候选位姿。利用路牌汉字笔画天然的横平竖直正交结构，将检测到的横笔画与竖笔画反投影到候选平面上，以正交性误差最小为准则从多组解中筛选出正确的位姿，无需额外标定或 IMU 辅助。}

\subsection{实战项目}
\cvline{自主路线规划数据采集无人机}{\textbf{项目简介：}基于树莓派与 RealSense 深度相机，在资源受限平台上二次开发 ORB-SLAM3 实现实时建图与位姿估计，结合 ROS 完成无人机全自动飞行与数据采集。主要工作：}
\cvline{}{{\textbf{硬件架构与系统集成：}树莓派 5 与 RealSense D435i 相机的软硬件调试；在 ARM 架构上完成 Ubuntu 系统移植，并自源码编译部署 ROS 及 RealSense 驱动，解决跨平台依赖冲突，实现传感器到计算芯片的低延迟数据链路。}}
\cvline{}{\textbf{SLAM 算法部署与建图：}针对树莓派 5 算力匮乏及易热降频的痛点，对 ORB - SLAM3 源码进行深度裁剪。主动关闭高耗能的全局回环检测线程，将单帧特征点提取上限由默认的 1000 裁剪至 400，并将局部 BA（光束法平差）的滑动优化窗口由 20 帧限制在 5 帧 以内。成功将算法平均解算帧率从严重的卡顿掉帧（3-5 fps）提升至稳定可用（12-15 fps），同时将 CPU 峰值占用率从 95\% 以上（极易触发温控降频）控制在 70\% 左右，以最低的算力代价满足了无人机低速飞行时对实时位姿反馈的及格线要求。}
\cvline{}{\textbf{ROS飞行控制与自动化采集：}编写 ROS 核心控制节点，订阅 ORB - SLAM3 发布的位姿话题（Topic），通过 MAVROS 实时向飞控发送位置控制指令（Setpoints），实现自主起飞与特定轨迹循迹；同步设计数据抓取脚本，实现无人机在飞行过程中全自动录制包含 RGB - D 图像、IMU 数据的 Rosbag 报文。}
% \cvline{}{\textbf{项目成果：}实现无人机全自动飞行数据采集任务，在端到端系统工程落地中积累了丰富的全栈开发经验。}

\cvline{基于多传感器的复杂环境三维建图系统}{\textbf{项目简介：}基于 3D 激光雷达、IMU 与视觉相机搭建多源传感器融合平台。主导底层的硬件组装、驱动编译、算法移植与代码重构，解决多传感器时空不同步与算力瓶颈问题，实现复杂工况下的高精度 3D 环境建图。}
\cvline{}{\textbf{多传感器硬件标定与时序同步：}负责底层硬件连线与驱动配置。独立使用标定工具完成雷达、六轴 IMU 与相机的空间外参标定；针对多源数据时间戳不一致问题，编写底层触发脚本，利用雷达 PPS（秒脉冲）信号触发相机，并解析 NMEA 报文打通硬件网络，将各传感器的时间戳误差压榨至微秒级（硬同步），消除高速运动下的点云重影。}
\cvline{}{\textbf{紧耦合 LIO（激光惯性里程计）算法开发：}将 FAST - LIO2 算法从开源库移植到实际计算平台并进行 C++ 源码级调试。针对测试中遇到的剧烈震动与长走廊（几何特征退化）场景，手动调整 IMU 噪声协方差矩阵参数，并修改体素滤波（Voxel Filter）的降采样叶子节点大小。通过调优 iEKF 紧耦合参数，成功解决设备在极端工况下的 Z 轴漂移问题。}
\cvline{}{\textbf{ROS 架构重构与零拷贝通信优化：} 主导底层系统提效，解决高线束雷达点云因 ROS 进程间通信（IPC）引发的严重卡顿。将原本独立的感知与建图 C++ 节点（Node）从底层代码改写为 ROS Nodelet 插件架构。利用 boost::shared\_ptr 共享内存指针，彻底省去庞大点云数据的网络序列化与反序列化环节（实现零拷贝传输），在建图精度无损的前提下，将系统整体 CPU 负载降低超 30\%。}

\cvline{基于多传感器的复杂环境三维建图系统}{\textbf{项目简介：}基于 3D 激光雷达、六轴 IMU 与工业相机搭建多源传感器融合平台。主导硬件集成、驱动编译、标定、算法适配与系统级优化，重点解决多传感器时间同步与实时算力瓶颈，实现复杂工况下高精度、低延迟的 3D 环境建图与动态障碍物滤除。}
\cvline{}{\textbf{多传感器硬件标定与时序同步：}负责传感器选型、连线与嵌入式驱动配置。利用 IMU 噪声标定工具校准六轴零偏，并基于标定板完成雷达、IMU 与相机的联合空间外参标定。针对多源数据时间基准不一致问题，采用 GPS 模块的 PPS（秒脉冲）信号与 NMEA 报文构建统一时钟源：PPS 信号同时接入雷达与相机的硬件触发/同步接口，NMEA 时间解析后写入设备时间戳缓冲，将传感器间时间同步误差控制在百微秒级以内。该硬同步方案为后续点云运动畸变校正及多模态融合提供了准确的时间基准。}
\cvline{}{\textbf{激光-惯性里程计适配与退化场景增强：}将 FAST-LIO2 算法从开源库移植到嵌入式计算平台，完成 C++ 源码编译、依赖适配与实时性调试。针对剧烈震动与长走廊等几何特征退化场景，手动校准 IMU 过程噪声协方差矩阵，并调整体素滤波叶子尺寸以改善点云匹配质量。为抑制 Z 轴累积漂移，在 iEKF 更新中引入地面平面约束因子，结合关键帧位姿图优化，将长距离走廊场景下的 Z 轴漂移控制在厘米级以内。}
\cvline{}{\textbf{ROS 系统架构重构与通信优化：}为解决高线束雷达点云因进程间通信（IPC）引发的严重卡顿，主导系统架构重构。将原本独立的雷达驱动、动态物体检测（视觉感知）与建图模块改写为 ROS Nodelet 插件架构，所有节点运行在同一进程内。利用 boost::shared\_ptr 实现节点间的点云零拷贝传递，消除序列化与反序列化开销。优化后，在 64 线激光雷达全量数据下，系统整体 CPU 占用率由 85\% 降至约 55\%，建图线程延迟降低 40\%，且建图精度不受影响。}

\subsection{实战项目}
\cvline{基于 ROS 2 的多模态感知与 3D 感知增强导航系统}{\textbf{项目简介：}基于树莓派 5 (Raspberry Pi 5 8GB) 边缘计算平台、RealSense D435i 深度相机与 RPLIDAR A2 单线激光雷达，构建 ROS 2 (Humble) 移动底盘。主导多源传感器融合、多模态 SLAM 部署与 Nav2 避障调优，利用深度相机补全单线雷达的感知盲区，实现在动态复杂场景下的稳定自主导航。主要工作：}

\cvline{}{\textbf{多源位姿融合与图优化 2D 建图：}基于 \texttt{robot\_localization} 框架，应用\textbf{扩展卡尔曼滤波（EKF）}融合底盘编码器与板载六轴 IMU，生成高频平滑的局部里程计（/odom），作为运动先验输入\textbf{Cartographer}。Cartographer 将激光帧与子图匹配，通过回环检测与 Ceres 优化器构建极低漂移的全局 2D 栅格地图（/map），并将 map→odom 变换单向发布修正 EKF 的累积漂移，避免闭环耦合。}

\cvline{}{\textbf{RGB-D 感知盲区修复（核心避障）：}针对 RPLIDAR A2 只能扫描固定高度平面、无法检测悬空物体（如桌角）的物理盲区，引入深度相机补全感知。基于 PCL 库对点云依次执行\textbf{PassThrough 直通滤波}（剔除超高空无用数据）与\textbf{VoxelGrid 体素降采样}，将精简后的 PointCloud2 注入 Nav2 体素层（Voxel Layer）。代价地图据此标记出原本被雷达漏检的障碍物投影（如桌角在膝盖高度形成的阻碍），使 DWA 规划器能提前避让，消除了纯 2D 雷达的感知盲区。}

\cvline{}{\textbf{全局导航与平滑避障控制：}重构 Nav2 行为树，上层采用\textbf{A* 算法}搜索全局最优路径，底层结合\textbf{DWA（动态窗口法）}执行局部轨迹采样与代价打分。通过精调 sim\_time 与代价函数权重，消除小车在动态人员干扰下的卡死现象。同时以低频异步模式（~1Hz）部署\textbf{RTAB-Map}，特征提取强制使用 ORB 以减少算力开销，仅用于生成带色彩纹理的三维点云供可视化展示，不参与避障决策。}
% \section{实习经历}
% \cventry{17.01-17.12}{北京九天探索}{组装和配置数据采集无人机}{}{}{工作中，
%   我负责了数据采集无人记得组装和配置，与采购部门对接，按照需求采购零件，组装检验无人机飞行稳定性。使用树莓派和realsense相机，编写自动飞行采集脚本，完成了无人机的自动飞行和数据采集任务。通过这段实习经历，我积累了丰富的无人机组装和配置经验，提升了我的动手能力和团队协作能力。}
% \cventry{16.01-16.12}{北京九天探索}{对工业相机进行标定和编写自动标定代码}{}{}{工作中，我按照开发部门的要求，对公司使用的工业相机进行标定，使用标定板进行拍摄，获取标定数据，并编写自动标定代码，实现了工业相机的自动标定功能。通过这段实习经历，我积累了丰富的工业相机标定经验，提升了我的编程能力和问题解决能力。}
\section{专业技能}
\cvline{\textbf{C++}}{熟练掌握 C++，具备丰富的项目实战经验和良好的面向对象编程习惯。}
\cvline{\textbf{ROS}}{熟悉 ROS 系统，能够使用 ROS 进行开发和调试，了解 ROS 的通信机制和工具。}
\cvline{\textbf{SLAM}}{熟悉 SLAM 基础理论，熟悉 ORB-SLAM3 的代码结构，能够在 ORB-SLAM3 的基础上进行二次开发。}
\cvline{\textbf{Linux Git}}{熟练使用 Linux 操作系统，熟悉常用的 Linux 命令和工具。熟练使用 Git 版本控制系统，能够进行代码的版本管理、分支操作和协作开发。}

\end{document}